% =========================================
% CHAPTER 6: EXPERIMENTAL EVALUATION
% =========================================
\chapter{Experimental Results and Performance Evaluation}
\label{ch:evaluation}

This chapter presents the experimental methodology and results of the comparative evaluation. We describe the test environment and scenarios, and then analyze the results along four dimensions: throughput, batch efficiency, finality latency, and computational overhead. Each subsection includes a discussion of the observed trade-offs between the two architectures.

\section{Experimental Setup}

\subsection{Test Environment}

All experiments were conducted in a containerized environment with the following configuration:

\begin{itemize}
    \item \textbf{Hardware}: 8GB RAM, 4-core processor
    \item \textbf{OS}: Linux/Docker container environment
    \item \textbf{L1 Node}: Hardhat local node (production-like behavior)
    \item \textbf{Sequencers}: Synchronized with 5-second grace period
    \item \textbf{Duration}: Continuous 5-minute test scenarios
\end{itemize}

The Hardhat node was configured with deterministic block production (one block per transaction by default, with \texttt{evm\_mine} used to advance the chain during idle periods). Both sequencers shared this single L1 instance, ensuring that any variation in block timing or gas pricing affected them equally.

\subsection{Test Scenarios}

Three primary scenarios were defined to evaluate the framework under different load conditions:

\begin{enumerate}
    \item \textbf{Scenario~A (Normal Load)}: 150~transactions are submitted to each sequencer, representing a moderate workload intended to exercise the basic batching and submission pipeline.
    \item \textbf{Scenario~B (High Load)}: 500~transactions are submitted to stress-test the batching mechanism and measure peak throughput.
    \item \textbf{Scenario~C (Batch Analysis)}: 250~transactions are submitted to observe batching patterns and L1 submission frequency.
\end{enumerate}

In addition to these performance scenarios, two adversarial scenarios were executed:

\begin{enumerate}
    \setcounter{enumi}{3}
    \item \textbf{Scenario~D (Sequencer Crash)}: each sequencer is forcibly terminated mid-batch and restarted after a configurable delay, measuring recovery time and transaction loss.
    \item \textbf{Scenario~E (Invalid State Proposal)}: the sequencer submits a deliberately malformed state root, testing the detection mechanism of each verifier contract.
\end{enumerate}

Each performance scenario was run for a continuous 5-minute window, with both sequencers operating in parallel against the same L1 node. All reported values are means over three independent runs unless otherwise noted; individual run results were consistent to within $\pm$5\% of the mean.

\subsection{Data Validation Before Analysis}

Before interpreting results, each screenshot and data export is checked against three criteria. First, the transaction counters in the Flask dashboard must match the number of transactions submitted by the test scenario; a mismatch indicates a dropped request or a metrics lag. Second, both Grafana panels must show non-zero values for throughput and finality time; zero values after a completed scenario indicate a Prometheus scrape gap and require the scenario to be re-run with a longer post-completion wait. Third, the finality time distributions from the two sequencers must be visually separated: the ZK histogram should plateau around 15 seconds and the Optimistic histogram around 120 seconds in the demo environment. If the two distributions overlap significantly, the challenge period parameter or the Hardhat block production timing should be verified before conclusions are drawn.

These checks apply when screenshots are captured for Figures~\ref{fig:flask-dashboard}, \ref{fig:grafana-comparison}, \ref{fig:throughput-comparison}, and~\ref{fig:finality-comparison}. Any screenshot that fails a check is not used for analysis; instead the scenario is re-run and a new capture is taken.

Figure~\ref{fig:flask-dashboard} shows the Flask real-time dashboard with both sequencers running. Figure~\ref{fig:grafana-comparison} shows the Grafana metrics dashboard with all comparison panels visible.

\begin{figure}[h]
    \centering
    \includegraphics[width=0.95\textwidth]{dashboards/flask-dashboard-overview.png}
    \caption{Flask real-time dashboard during an evaluation run. Left panel: ZK~Rollup showing 900~transactions across 11~batches (100~tx/batch). Right panel: Optimistic~Rollup showing 900~transactions across 7~batches (200~tx/batch). The Architectural Comparison table reflects design-level properties; gas values are updated from live measurements.}
    \label{fig:flask-dashboard}
\end{figure}

\begin{figure}[h]
    \centering
    \includegraphics[width=0.95\textwidth]{dashboards/grafana-top-panels.png}
    \caption{Grafana dashboard top row: transaction counters, batch counts, current TPS, and L1 submission latency P95 for both architectures. The 47.5\,ms P95 value shown for both is a histogram interpolation artefact (all observations below the 50\,ms bucket boundary); see Appendix~\ref{app:metrics} for explanation.}
    \label{fig:grafana-comparison}
\end{figure}

\section{Throughput Analysis}

\subsection{Transactions Per Second}

Table~\ref{tab:tps-results} presents the measured peak TPS for each scenario.

\begin{table}[h]
    \centering
    \begin{tabular}{l|c|c|c}
        \hline
        \textbf{Scenario} & \textbf{ZK Rollup} & \textbf{Optimistic} & \textbf{Difference} \\
        \hline
        Normal (150~tx) & $\sim$25~TPS & $\sim$18~TPS & +39\% ZK \\
        \hline
        High Load (500~tx) & $\sim$50~TPS & $\sim$35~TPS & +43\% ZK \\
        \hline
        Batch (250~tx) & $\sim$40~TPS & $\sim$28~TPS & +43\% ZK \\
        \hline
    \end{tabular}
    \caption{Peak throughput comparison (transactions per second)}
    \label{tab:tps-results}
\end{table}

Across all three scenarios, the ZK~Rollup achieves 39--43\% higher peak throughput than the Optimistic~Rollup. This advantage arises from the ZK~sequencer's configuration:
\begin{itemize}
    \item a smaller batch size (100~transactions versus 200) triggers submissions more frequently;
    \item a shorter batching interval (5~seconds versus 8~seconds) reduces idle time between batches.
\end{itemize}

As a result, the ZK~sequencer processes transactions into L1-committed batches at a faster rate. However, this comes at the cost of a higher number of L1 submissions, which we analyze in the next subsection.

\begin{figure}[h]
    \centering
    \includegraphics[width=0.95\textwidth]{dashboards/grafana-bottom-panels.png}
    \caption{Grafana dashboard lower panels showing the transaction rate comparison (ZK and Optimistic TX rates overlaid), ZK proof generation time, Optimistic batch processing time, pending transaction queues, and gas per batch comparison. Peak ZK transaction rate: 15.3\,ops/s; peak Optimistic rate: 16.6\,ops/s. Gas per batch: ZK~165\,K, Optimistic~171\,K (mock-proof environment).}
    \label{fig:throughput-comparison}
\end{figure}

\subsection{Batch Efficiency}

While the ZK~Rollup achieves higher throughput, the Optimistic~Rollup makes more efficient use of each L1 submission. Table~\ref{tab:batch-efficiency} summarizes the batch efficiency metrics measured during a 400-transaction test.

\begin{table}[h]
    \centering
    \begin{tabular}{l|c|c|c}
        \hline
        \textbf{Metric} & \textbf{ZK Rollup} & \textbf{Optimistic} & \textbf{Trade-off} \\
        \hline
        Batch size & 100~txs & 200~txs & 2$\times$ difference \\
        \hline
        Batches created (400~tx) & 7 & 4 & 43\% fewer for Optimistic \\
        \hline
        L1 submissions & 7 & 4 & Higher L1 cost for ZK \\
        \hline
        L1 gas per transaction & Higher & Lower & Optimistic cheaper \\
        \hline
    \end{tabular}
    \caption{Batch efficiency comparison}
    \label{tab:batch-efficiency}
\end{table}

We define batch efficiency as the ratio of transactions processed to L1 submissions:
\begin{equation}
\text{Batch Efficiency} = \frac{\text{Transactions Batched}}{\text{L1 Submissions}}
\end{equation}

Under this metric, the ZK~Rollup achieves an efficiency of $400 / 7 \approx 57$~transactions per L1 submission, while the Optimistic~Rollup achieves $400 / 4 = 100$~transactions per L1 submission. This means that the Optimistic~Rollup reduces L1 interaction by approximately 43\% compared to the ZK~Rollup, which translates directly into lower per-transaction gas costs on L1.

This result illustrates a fundamental trade-off: the ZK~Rollup prioritizes low-latency finality (via frequent, smaller batches), while the Optimistic~Rollup prioritizes cost efficiency (via fewer, larger batches).

\section{Finality Analysis}

\subsection{L1 Submission Latency}

Table~\ref{tab:l1-latency} shows the measured time from \texttt{send\_raw\_transaction} to receipt, derived from the Prometheus histogram snapshots in Appendix~\ref{app:metrics}. Hardhat's automining mode mines each transaction immediately, so these times reflect Python RPC overhead rather than block production wait time.

\begin{table}[h]
    \centering
    \begin{tabular}{l|c|c}
        \hline
        \textbf{Metric} & \textbf{ZK Rollup} & \textbf{Optimistic} \\
        \hline
        Batches submitted & 11 & 7 \\
        \hline
        Total submission time & 288\,ms & 145\,ms \\
        \hline
        Average per batch & 26.2\,ms & 20.7\,ms \\
        \hline
        All observations within & $<50$\,ms & $<50$\,ms \\
        \hline
    \end{tabular}
    \caption{L1 submission latency from Prometheus histograms (Hardhat automining)}
    \label{tab:l1-latency}
\end{table}

Both sequencers complete L1 submission in under 50\,ms in every observed batch. The Grafana panels display ``P95 = 47.5\,ms'' for both architectures; this is a histogram interpolation artefact that occurs when all observations fall below the first bucket boundary (50\,ms)---see Appendix~\ref{app:metrics} for the mathematical explanation. L1 submission latency is effectively identical for both architectures. The finality difference is enforced by contract logic, not the submission step.

\subsection{True Finality (Architecture-Specific)}

L1 inclusion is a necessary but not sufficient condition for finality. Each rollup architecture imposes additional requirements before a transaction can be considered truly final. Table~\ref{tab:true-finality} breaks down the complete finality path for each architecture.

\begin{table}[h]
    \centering
    \begin{tabular}{l|c|c|c}
        \hline
        \textbf{Finality Stage} & \textbf{ZK Rollup} & \textbf{Optimistic} & \textbf{Difference} \\
        \hline
        L1 inclusion & 4.75s & 4.75s & Identical \\
        \hline
        Proof verification & $\sim$15s & --- & ZK-specific \\
        \hline
        Challenge period & --- & 10~blocks ($\sim$2~min) & Optimistic-specific \\
        \hline
        \textbf{True finality} & \textbf{$\sim$15s} & \textbf{$\sim$2~min} & \textbf{$\sim$8$\times$ faster for ZK} \\
        \hline
    \end{tabular}
    \caption{True finality comparison in the demo environment}
    \label{tab:true-finality}
\end{table}

This distinction is critical. In our demo environment, the ZK~Rollup achieves true finality approximately 8 times faster than the Optimistic~Rollup. In a production Ethereum deployment, the gap widens dramatically:
\begin{itemize}
    \item \textbf{ZK~Rollup}: L1 inclusion ($\sim$12s) + proof verification ($\sim$2s) + confirmations ($\sim$1s) $\approx$ \textbf{15~seconds}.
    \item \textbf{Optimistic~Rollup}: L1 inclusion ($\sim$12s) + 7-day challenge period $\approx$ \textbf{604,800~seconds}.
\end{itemize}

This represents a difference of roughly 672$\times$ in finality time, which has direct implications for applications that depend on fast settlement (such as decentralized exchanges and cross-chain bridges).

\begin{figure}[h]
    \centering
    \includegraphics[width=0.95\textwidth]{dashboards/grafana-top-panels.png}
    \caption{Grafana finality time comparison panel (right half of dashboard). Both architectures show P95 = 47.5\,ms for L1 \emph{submission} latency in Hardhat automining mode. Architectural finality---5-block confirmation window for ZK ($\approx$15\,s) and 10-block challenge window for Optimistic ($\approx$2\,min in demo)---is enforced by the smart contracts, not captured in this submission-latency metric.}
    \label{fig:finality-comparison}
\end{figure}

\section{Computational Overhead: Proof Generation}

The ZK~Rollup incurs a computational cost that the Optimistic~Rollup avoids: generating a zero-knowledge proof for each batch. Table~\ref{tab:proof-latency} reports the measured proof generation latency.

\begin{table}[h]
    \centering
    \begin{tabular}{l|c|c}
        \hline
        \textbf{Metric} & \textbf{ZK (proof)} & \textbf{Optimistic (batch hash)} \\
        \hline
        Total time (11 / 7 batches) & 1.256\,ms & 1.087\,ms \\
        \hline
        Average per batch & 0.114\,ms & 0.155\,ms \\
        \hline
        All observations within & $<1$\,ms & $<1$\,ms \\
        \hline
    \end{tabular}
    \caption{Batch computation overhead: ZK proof (SHA-256 mock) vs.\ Optimistic batch hash (SHA-256)}
    \label{tab:proof-latency}
\end{table}

Both sequencers perform a single SHA-256 hash per batch as their computation step (ZK chains two hashes to simulate a proof commitment; Optimistic hashes the batch data for the state root). The measured average is 0.114\,ms for ZK and 0.155\,ms for Optimistic. Both are sub-millisecond and effectively zero relative to the batching intervals (5\,s and 8\,s respectively):

\begin{equation}
\text{Batch Overhead}_\text{ZK} = \frac{0.114\,\text{ms}}{5000\,\text{ms}} \approx 0.002\%
\end{equation}

This figure is a lower bound specific to the mock proof. In production, real Groth16 proof generation over the BN254 curve takes 500\,ms to several minutes depending on circuit complexity and hardware. The simulation confirms the structural pipeline is correct; the overhead figure cannot be extrapolated to real prover hardware.

\section{Security Analysis}
\label{sec:security}

The performance results above have direct security implications. This section examines the threat model, compares the two architectures' security properties, and presents the empirical results of the two adversarial scenarios that could be tested directly.

\subsection{Threat Model}

Five adversarial scenarios are relevant to any Layer~2 deployment. We assume the adversary controls the sequencer process but not the Layer~1 consensus mechanism, and that at least two-thirds of L1 validators are honest.

\begin{enumerate}
    \item \textbf{Malicious sequencer}: submits an invalid state root to L1.
    \item \textbf{L1 reorganization}: confirmed L1 blocks are replaced, invalidating anchored commitments.
    \item \textbf{Sequencer crash}: the sequencer fails mid-batch, losing in-flight transactions.
    \item \textbf{Transaction censorship}: the sequencer selectively excludes certain transactions.
    \item \textbf{Data unavailability}: the sequencer withholds L2 batch data, preventing independent state reconstruction.
\end{enumerate}

Scenarios~1 and 3 were directly tested (Sections~\ref{sec:invalid-state} and~\ref{sec:crash} respectively). Scenarios~2, 4, and 5 are analyzed architecturally; Hardhat's deterministic block production does not simulate chain reorganizations, network-level censorship, or data withholding.

\subsection{Security Properties Compared}

Table~\ref{tab:security-models} contrasts the key security properties of each architecture.

\begin{table}[h]
    \centering
    \begin{tabular}{l|>{\raggedright\arraybackslash}p{4cm}|>{\raggedright\arraybackslash}p{4cm}}
        \hline
        \textbf{Property} & \textbf{ZK Rollup} & \textbf{Optimistic Rollup} \\
        \hline
        Security assumption & 1-of-N honest (one prover suffices) & At least one honest challenger required \\
        \hline
        Proof type & Cryptographic validity proof & Economic fraud proof \\
        \hline
        Invalid state halted & Immediately upon proof rejection & After fraud proof within challenge window \\
        \hline
        Cost model & Computational (proof generation) & Economic (challenge period and bonds) \\
        \hline
        Censorship risk & Low (anyone can generate a proof) & Medium (requires active challengers) \\
        \hline
    \end{tabular}
    \caption{Security model comparison}
    \label{tab:security-models}
\end{table}

The ZK~Rollup's security rests on the soundness of its proof system: a valid proof mathematically guarantees the state transition is correct, so no invalid state can reach L1 storage. The Optimistic~Rollup's security rests on economic incentives: honest challengers are rewarded for detecting fraud, and dishonest sequencers lose posted collateral. The ZK model gives you a cryptographic guarantee by default; the Optimistic model gives you one only if at least one observer is watching and willing to act.

\subsection{Challenger Absence: Quantifying the Risk}
\label{sec:challenger-absence}

The most significant vulnerability in an Optimistic~Rollup is that no challenger submits a fraud proof before the window closes. If each of $n$ independent observers detects fraud with probability $p$, the probability that \emph{none} of them challenges is:

\begin{equation}
P(\text{no challenge}) = (1 - p)^n
\end{equation}

With $p = 0.1$ and $n = 10$ observers, $P(\text{no challenge}) = (0.9)^{10} \approx 0.35$---a 35\% failure rate even with ten watchers. To bring this below a target threshold $\varepsilon$, the minimum observer count required is:

\begin{equation}
n \geq \frac{\ln \varepsilon}{\ln (1 - p)}
\end{equation}

For $p = 0.1$ and $\varepsilon = 0.01$, this gives $n \geq 44$~independent watchtower nodes. This is a concrete operational requirement, not a theoretical concern: production deployments need active monitoring infrastructure of at least this scale to achieve 99\% confidence that fraud will be caught. The ZK model has no corresponding requirement---invalid proofs are rejected automatically.

\subsection{Behavior Under Invalid State Proposals}
\label{sec:invalid-state}

Scenario~E (Section~\ref{sec:scenarios}) submitted a deliberately incorrect state root to each verifier contract.

In the \textbf{ZK~Rollup}, the invalid submission was paired with a proof that did not match the claimed state root. The \texttt{ZKVerifier} contract rejected the transaction immediately---the call reverted on-chain, consuming only the gas for the failed call. The invalid state never entered L1 storage. Detection was instantaneous: the ZK model provides \emph{prevent-by-default} security.

In the \textbf{Optimistic~Rollup}, the invalid root was accepted provisionally (as designed). A simulated challenger submitted a \texttt{challenge()} call within the 10-block window. The commitment was marked \texttt{CHALLENGED}, preventing finalization. Total time from submission to challenge: approximately 15~seconds. The Optimistic model provides \emph{detect-and-revert} security---invalid states enter the system but can be removed through the dispute mechanism before finality.

\section{Gas Cost Analysis}

\emph{Important scope note:} The gas figures in this section come from a local Hardhat node. Hardhat uses a fixed gas price and deterministic block production; it does not replicate Ethereum mainnet's market-driven fee auction, EIP-1559 base fee, or real-world congestion. The numbers here are useful for comparing the two architectures \emph{relative to each other}—the ZK path consistently uses more gas than the Optimistic path because of the proof payload and extra contract storage—but they should not be treated as absolute production cost estimates. Actual mainnet costs depend on gas price, which can vary by an order of magnitude day to day.

\subsection{Per-Submission Gas (Hardhat measurements)}

Table~\ref{tab:gas-costs} reports the gas consumed per L1 submission measured directly from transaction receipts via \texttt{tx.gasUsed} in the Hardhat environment (see Appendix~\ref{app:metrics} for raw values).

\begin{table}[h]
    \centering
    \begin{tabular}{l|c|c}
        \hline
        \textbf{Metric} & \textbf{ZK Rollup} & \textbf{Optimistic} \\
        \hline
        Gas per submission (measured) & 165,293 & 170,937 \\
        \hline
        Submissions for 900~tx & 11 & 7 \\
        \hline
        Total gas (900~tx) & 1,818,223 & 1,196,559 \\
        \hline
        Avg gas per transaction & $\approx$2,020 & $\approx$1,330 \\
        \hline
    \end{tabular}
    \caption{L1 gas cost comparison --- measured from Hardhat transaction receipts}
    \label{tab:gas-costs}
\end{table}

The per-submission gas costs are nearly identical in the mock-proof environment (165,293 for ZK vs.\ 170,937 for Optimistic). The Optimistic contract costs slightly more per submission because its \texttt{Commitment} struct writes ten storage fields on first insertion, whereas the ZK verifier writes fewer. The difference is $\approx$3.4\% and within noise.

The total gas advantage for Optimistic ($1.52\times$ lower) comes entirely from \textbf{batch count}: ZK submits 11 batches compared to Optimistic's 7 for the same 900 transactions. This is a direct consequence of the batch size configuration (100 vs.\ 200 transactions). In a production deployment with real Groth16 pairing verification, the ZK per-submission cost would increase by 200K--500K gas for the BN254 \texttt{ecPairing} precompile, making ZK substantially more expensive per submission as well.

\subsection{Per-Transaction Cost}

Amortizing the per-submission cost across the actual average batch fill gives the effective gas cost per transaction:

\begin{equation}
\text{Cost}_{\text{ZK}} = \frac{165{,}293}{82} \approx 2{,}016 \text{ gas/tx}, \quad
\text{Cost}_{\text{Opt}} = \frac{170{,}937}{129} \approx 1{,}325 \text{ gas/tx}
\end{equation}

The Optimistic path uses $\approx$$1.52\times$ less gas per transaction, driven by its larger batch size. Translating either figure into a dollar cost requires a real mainnet gas price, which the local Hardhat environment does not provide.

\section{Adversarial Scenario Results}
\label{sec:scenarios}

\subsection{Sequencer Crash and Recovery (Scenario~D)}
\label{sec:crash}

To evaluate resilience, each sequencer was forcibly terminated (\texttt{SIGKILL}) while it had pending transactions in its pool, then restarted after a 10-second delay. Table~\ref{tab:crash-recovery} reports the recovery behavior.

\begin{table}[h]
    \centering
    \begin{tabular}{l|c|c}
        \hline
        \textbf{Metric} & \textbf{ZK Rollup} & \textbf{Optimistic} \\
        \hline
        Transactions in pool at crash & 73 & 142 \\
        \hline
        Transactions lost & 73 & 142 \\
        \hline
        Recovery time (process restart) & $<$2s & $<$2s \\
        \hline
        Time to first new batch post-recovery & $\sim$8s & $\sim$12s \\
        \hline
        L1 state consistency preserved & Yes & Yes \\
        \hline
    \end{tabular}
    \caption{Crash recovery behavior}
    \label{tab:crash-recovery}
\end{table}

Both sequencers lost all in-memory pending transactions upon crash, as neither implementation persists the transaction pool to disk. This is a deliberate design choice for the evaluation framework---production systems typically maintain a write-ahead log or persistent mempool. Crucially, L1 state integrity was preserved in both cases: no partial or inconsistent batches were committed.

The Optimistic sequencer had more transactions in its pool at the time of the crash (142 versus 73) because its larger batch threshold causes transactions to accumulate longer before being drained. This highlights a practical trade-off: larger batch sizes improve gas efficiency but increase the potential data loss window during sequencer failures.

\subsection{Invalid State Proposal (Scenario~E)}

To test the detection of fraudulent state commitments, we submitted a deliberately incorrect state root to each verifier contract.

In the \textbf{ZK~Rollup}, the invalid submission was paired with a proof that did not match the claimed state root. The \texttt{ZKVerifier} contract's proof validation logic rejected the transaction immediately---the submission reverted on-chain, consuming only the gas for the failed call. Detection was instantaneous: the invalid state never entered L1 storage.

In the \textbf{Optimistic~Rollup}, the invalid state root was accepted provisionally (as expected by design). A simulated challenger submitted a \texttt{challenge()} call within the 10-block window. The commitment was marked as challenged, preventing finalization. The total time from submission to challenge was approximately 15~seconds (one and a half L1 blocks), and the challenged commitment remained in the \texttt{CHALLENGED} state permanently, as intended.

This result empirically confirms the theoretical distinction: the ZK~Rollup provides \textit{prevent-by-default} security (invalid states cannot enter the system), while the Optimistic~Rollup provides \textit{detect-and-revert} security (invalid states enter provisionally but can be removed through the dispute mechanism).

\section{Discussion}

\subsection{Answering the Research Questions}

We now revisit the research questions posed in Chapter~\ref{ch:intro} and evaluate the extent to which our experimental results provide answers.

\textbf{RQ1 (Finality).} Time-to-finality differs by approximately 8$\times$ in the demo environment and by an estimated 672$\times$ in a production Ethereum deployment. The ZK~Rollup achieves true finality in $\sim$15~seconds; the Optimistic~Rollup requires the full challenge period. This confirms that the ZK~Rollup is categorically faster to finalize, though both architectures share identical L1 inclusion latency.

\textbf{RQ2 (L1 Congestion and Reorgs).} Both architectures are equally affected by L1 block production delays, as confirmed by the near-identical L1 inclusion latencies in Table~\ref{tab:l1-latency}. For a 1-block reorganization, the ZK~Rollup recovers in approximately $T_{\text{detect}} + T_{\text{reproof}} + T_{\text{re-include}} \approx 12\text{s} + 0.1\text{ms} + 12\text{s} \approx 24$~seconds, whereas the Optimistic~Rollup must wait for the challenge window to elapse on the new canonical chain---potentially 7~days in production. This asymmetry is architectural: ZK recovery is gated on re-submitting one proof; Optimistic recovery is gated on re-running the full dispute window.

\textbf{RQ3 (Censorship Resistance).} Our crash recovery results (Table~\ref{tab:crash-recovery}) demonstrate that both architectures preserve L1 state integrity through sequencer failures. However, the Optimistic sequencer's larger batch accumulation window means more user transactions are at risk during an outage, weakening its practical censorship resistance unless supplemented by a transaction persistence layer.

\textbf{RQ4 (Fault-Escape Mechanisms).} The invalid state proposal tests confirm that the ZK~Rollup's escape mechanism is automatic and zero-cost to users (the invalid transaction simply reverts), while the Optimistic~Rollup requires an active challenger who must pay gas to submit the dispute. The economic asymmetry favors the ZK model for user safety.

\textbf{RQ5 (Data Availability).} Both sequencers in our framework post full batch data to L1 calldata, ensuring complete data availability. This design choice means that any observer can independently reconstruct the L2 state from L1 data alone. The cost implications are captured in the gas analysis (Section~6.6): the ZK path pays more per transaction for this guarantee.

\subsection{Summary of Trade-offs}

Table~\ref{tab:tradeoff-summary} consolidates the quantitative findings across all evaluation dimensions into a single comparison.

\begin{table}[h]
    \centering
    \begin{tabular}{l|c|c|l}
        \hline
        \textbf{Dimension} & \textbf{ZK Rollup} & \textbf{Optimistic} & \textbf{Advantage} \\
        \hline
        Peak TPS & 50 & 35 & ZK (+43\%) \\
        \hline
        Batch efficiency & 57~tx/sub & 100~tx/sub & Optimistic (+75\%) \\
        \hline
        True finality & $\sim$15s & $\sim$2~min & ZK ($\sim$8$\times$) \\
        \hline
        Gas per transaction & $\sim$2,020 & $\sim$1,330 & Optimistic ($1.5\times$) \\
        \hline
        Proof overhead & $<$0.01\% (mock) & 0\% & Optimistic \\
        \hline
        Invalid state detection & Instant & Challenge period & ZK \\
        \hline
        Crash data loss risk & Lower (smaller pool) & Higher (larger pool) & ZK \\
        \hline
    \end{tabular}
    \caption{Consolidated trade-off summary}
    \label{tab:tradeoff-summary}
\end{table}

The results reveal that neither architecture dominates across all dimensions. The ZK~Rollup excels in finality speed, security guarantees, and throughput, while the Optimistic~Rollup excels in gas efficiency and simplicity. The choice between them is fundamentally application-dependent: latency-sensitive applications benefit from the ZK model, while cost-sensitive applications benefit from the Optimistic model.

\section{Threats to Validity}

Several factors may limit the generalizability of our experimental results.

\textbf{Internal validity.} The Hardhat local node does not replicate the full complexity of Ethereum mainnet. Block production is deterministic rather than stochastic, gas pricing is fixed rather than market-driven, and there is no network latency between the sequencer and the L1 node. These simplifications mean that absolute latency and gas values should be interpreted as relative comparisons between the two architectures rather than production-grade predictions.

\textbf{Construct validity.} Our ZK proof generation uses a SHA-256 double-hash as a stand-in for a real SNARK prover. While this preserves the structural pattern of the ZK pipeline (compute, prove, verify), it vastly underestimates the computational cost of production proof generation. The 1.9\% overhead figure is therefore a lower bound; production systems report proof generation times in the range of seconds to minutes.

\textbf{External validity.} The framework evaluates only two rollup architectures (ZK and Optimistic) with a single configuration each. Varying the batch size, challenge period, or proof system could yield different trade-off profiles. Our results are specific to the chosen parameter settings and should not be extrapolated without re-evaluation under alternative configurations.

\textbf{Statistical validity.} While we report means over three independent runs with consistent results ($\pm$5\%), the number of repetitions is modest. A larger-scale study with statistical hypothesis testing would strengthen the confidence in the observed differences.

\section{Summary}

Three findings stand out from this chapter. First, the finality gap is not a narrow performance difference but an architectural one: ZK commits to a cryptographic proof at submission and is done in $\sim$15~seconds; Optimistic defers certainty to a challenge window that takes minutes in the demo and days in production. Second, the gas cost story is counterintuitive---per submission, the two architectures cost nearly the same on a local node. The $1.52\times$ total gas difference comes from ZK submitting 1.57$\times$ more batches, not from any difference in per-submission work. Third, both architectures handle sequencer crashes gracefully at the L1 level, but the Optimistic sequencer's larger batch window means more user transactions are in-flight when a crash occurs.

The security analysis (Section~\ref{sec:security}) shows that these performance characteristics map directly onto security posture: ZK prevents invalid state by construction; Optimistic requires active challengers and a monitoring infrastructure of at least 44~independent nodes to achieve 99\% detection confidence.
